\documentclass[11pt,a4paper]{article}
\usepackage[T1]{fontenc}
\usepackage{lmodern}
\usepackage[margin=23mm]{geometry}
\usepackage{amsmath,amssymb,booktabs,microtype}
\usepackage[hidelinks,hypertexnames=false]{hyperref}
\pdfinfoomitdate=1
\pdftrailerid{}
\pdfsuppressptexinfo=15
\title{Cloudy comparison: from a discrepancy to a physical repair}
\author{exojax-escape}
\date{21 November 2026}
\begin{document}
\maketitle

\section{What is being tested}
The first finite hydrogen atom solved its declared equations, but underpredicted
Cloudy's $n=3$ population by about 30\% in a controlled optically thin case.
A converged solver and an agreeing physical model are different claims. This
investigation asks which rate choices cause that difference, repairs the
identified choices using published equations, and repeats a fixed comparison.
It does not make Cloudy experimental truth or establish a planetary atmosphere.

The comparison holds temperature, hydrogen density and diluted blackbody
radiation fixed. The baseline is $T=9000$ K, $n_{\rm H}=10^{12}$ cm$^{-3}$,
$T_{\rm rad}=10000$ K and $W=0.025$, at maximum shells 4 and 5.
Six further tests vary one coordinate: $T=6000,12000$ K,
$n_{\rm H}=10^{10},10^{14}$ cm$^{-3}$, or $W=0.01,0.05$ at maximum shell 5.
These eight cases and a 5\% tolerance for both $n_2$ and $n_3$ were recorded
before measuring the repaired validation results. They are eight samples,
not a validated Cartesian product of parameter ranges. Each reference also
needs a matched frequency refinement changing its populations by less than 1\%.

\section{Following the causal chain}
First, the actual final Cloudy rate operator and source vector are extracted.
An independent linear solve of those same quantities reproduces Cloudy's
populations to about $1.2\times10^{-13}$ relative error. This isolates the
population solver from differences in atomic coefficients.

Next, one physical rate block is exchanged at a time. Both directions of a
collision pair are changed together; photoionization and its inverse
recombination are changed together. The experiment is repeated at fixed
actual electron density and with charge closure. In this diagnostic charge
closure, the extracted coefficients, including their screening state, remain
frozen; it is not a fresh Cloudy rate evaluation at each trial density. A reverse intervention
removes the same block from the reference. This is more discriminating than
changing many parameters until a population matches.

At the $n_{\max}=5$ baseline, adding only the omitted $n=3$ orbital mixing explains about
88\% of the original absolute $n_3$ discrepancy with charge closure; removing
it from Cloudy restores about 90\%. The fixed-electron-density tests give
about 83\% and 86\%. The predeclared 80\% dominance criterion is therefore met
\emph{at this baseline}. Replacing the bound-free block as well reduces the
remaining discrepancy to about 0.28\% in this diagnostic experiment.
Those block-swap coefficients are recorded diagnostic quantities; they are
not used as a fitted runtime table.

The first held-out campaign exposes a second failure: in the dilute case,
$n_2$ is about 12.46\% too high while $n_3$ agrees. Swapping only the paired
$n=2$ mixing block removes about 96.8\% of that discrepancy, and the reverse
substitution restores about 96.3\%. The legacy Huang/Seaton electron-plus-proton
fit \cite{huang} and the selected Cloudy PSM20 proton closure are different physical
approximations. This is a model-choice discrepancy, not an incorrect
implementation of the old published fit. The second experiment was declared
before measurement in the agent log; its JSON follow-up record was written
after its outcome and is explicitly labeled accordingly.

\begin{figure}[ht]
\centering
\begin{tabular}{c@{\quad$\longrightarrow$\quad}c}
\fbox{\parbox{.36\textwidth}{\centering Retained failed comparison\\$n_3$ about 30\% low}}
&\fbox{\parbox{.36\textwidth}{\centering Isolate one rate block\\Forward and reverse experiments}}\\[5mm]
\multicolumn{2}{c}{$\downarrow$}\\[3mm]
\fbox{\parbox{.36\textwidth}{\centering Adopt a sourced equation\\PSM20 mixing and Coulomb dipoles}}
&\fbox{\parbox{.36\textwidth}{\centering Repeat unchanged cases\\A new $n_2$ failure reopens diagnosis}}\\[5mm]
\multicolumn{2}{c}{$\downarrow$}\\[3mm]
\multicolumn{2}{c}{\fbox{\parbox{.80\textwidth}{\centering Numerical checks + all cases + reference refinement\\Only then review the scoped agreement claim}}}
\end{tabular}
\caption{The investigation loop. A failed held-out case remains evidence and
selects the next experiment; it does not end the task or disappear from the case list.}
\end{figure}

\section{What changed in the physical equations}
\subsection{Orbital mixing connects different lifetimes}
The $3s$, $3p$ and $3d$ states can radiate through different decay paths.
Proton collisions move atoms between these states and thereby alter how long
atoms remain in the combined $n=3$ reservoir. Redistribution can therefore
change its total population, even when collisions themselves conserve atoms.
The same reasoning matters for $2s$ and $2p$, whose radiative lifetimes differ
strongly.

The new finite atom uses the PSM20 impact-parameter approximation
\cite{badnell} for the three pairs $2p\leftrightarrow2s$,
$3p\leftrightarrow3s$ and $3d\leftrightarrow3p$. It averages a specified
collision cross section over Maxwellian speeds, with maximum impact parameter
limited by both electron Debye screening and the departing state's lifetime:
\[
R_c(E)=\min[R_D(T,n_e),R_\tau(E)],\qquad
q(T,n_e)=\sqrt{\frac{8}{\pi\mu}}(k_{\rm B}T)^{-3/2}
\int_0^\infty E\sigma(E)e^{-E/(k_{\rm B}T)}\,dE.
\]
The full cross section is in the generated appendix. This is an approximate
near-degenerate atomic model, particularly for low principal quantum number;
matching Cloudy's choice does not prove exact low-$n$ quantum scattering.
The older six-state API retains its documented Huang/Seaton fit.

Cloudy C23 switches to a pure-Debye shortcut in one part of this calculation.
Near its switching point that produces a rate jump of about 6.6\% for one
$n=3$ transition. The differentiable implementation integrates the combined
cutoff at every density, avoiding that numerical shortcut. It also uses the
neutral-protium/proton reduced mass rather than standard atomic weight.
These small declared differences are checked separately from the code that
reproduces Cloudy's exact convention. We do not copy an artificial discontinuity
merely to obtain identical output.

\subsection{Ionization and recombination share orbital-resolved dipoles}
The preceding Kramers approximation treated each principal shell as having
one photoionization cross section. The new default evaluates hydrogenic
Coulomb dipoles with the Brocklehurst radial recurrence
\cite{brocklehurst,cloudy}. It resolves the first six orbital states and
statistically averages the collapsed $n=4,5$ shells:
\[
\overline\sigma_n(\nu)=\sum_{\ell=0}^{n-1}
\frac{2\ell+1}{n^2}\sigma_{n\ell}(\nu).
\]
The same $\sigma_i(\nu)$ supplies photoionization, inverse Milne recombination,
continuum absorption and recombination emission. Changing only one of these
would break the shared rate/transfer convention. The old Kramers option remains
available for controlled comparisons. Neither version rescales populations
to match the reference.

\subsection{Density dependence belongs inside differentiation}
At a trial electron density, the bound-state equations are linear in
$b_i=n_i/n_e$. Nuclear conservation leaves a scalar root for
$z=\log(n_e/n_{\rm H})$:
\[
F(z;\theta)=z+\log\left[1+\sum_i b_i(n_{\rm H}e^z;\theta)\right]=0,
\qquad \frac{\partial z}{\partial\theta}=-\frac{F_\theta}{F_z}.
\]
The screened mixing coefficients must be recalculated at the trial $n_e$ both
in the residual and in its derivative. Freezing them at the final density can
reproduce the value while giving the wrong gradient. The canonical population
API is tested against finite differences and against that deliberately frozen
control. The coupled model uses the same kernels. This supports later gradient
inference; it does not by itself demonstrate NUTS convergence.

An execution-only follow-up changes how the exponential integral is batched.
Ordinary vectorization evaluates expensive inactive special-function branches
on this CPU/JAX configuration. Sequential scalar batching underneath the
existing analytic derivative preserves the equations and reduces the measured
three-element kernel time from about 8.42 ms to 37.6 microseconds; this is not
an end-to-end sampler timing. Representative first and second derivatives,
canonical population derivatives and fresh eight-case outputs are checked.
All populations and electron densities change by at most
$2.20\times10^{-14}$ relative in that replay. The exploratory performance
measurements preceded the plan author receiving their results; the retained
chronology correction explicitly withdraws an initial pre-outcome claim.
Old and optimized source identities have separate Subjects and Bindings.

\section{Repeated comparison and the remaining differences}
The final source-pinned comparison passes the original 5\% population criterion
in all eight cases. The maximum error is 1.100\%; matched Cloudy frequency
refinement changes any population by at most 0.049\%.
\begin{center}
\begin{tabular}{lrr}
\toprule
Case & $|\Delta n_2|/n_{2,\mathrm{ref}}$ (\%) & $|\Delta n_3|/n_{3,\mathrm{ref}}$ (\%)\\
\midrule
Baseline, $n_{\max}=4$ & 0.322 & 0.364 \\
Baseline, $n_{\max}=5$ & 0.446 & 0.508 \\
$T=6000$ K & 0.422 & 0.322 \\
$T=12000$ K & 0.983 & 1.100 \\
$n_{\rm H}=10^{10}$ cm$^{-3}$ & 0.436 & 0.376 \\
$n_{\rm H}=10^{14}$ cm$^{-3}$ & 0.186 & 0.754 \\
$W=0.01$ & 0.315 & 0.367 \\
$W=0.05$ & 0.523 & 0.588 \\
\bottomrule
\end{tabular}
\end{center}
These results establish agreement for the recorded finite-atom samples only.
A final baseline rate-budget experiment finds that exchanging the bound-free
block changes the $n_3$ residual by about 0.28 percentage points and exchanging
nonmixing bound-bound collisions changes it by about 0.17 points. The main
remaining choices are Cloudy's ground-state Verner fit versus the Coulomb
cross section, and added-shell collision data versus the shell/branch
approximation. The one-at-a-time shifts are not additive. Some individual
added-shell collision coefficients still differ by 76--95\%, even while the
summed populations meet the declared criterion. Thus agreement of $n_2,n_3$
is explicitly narrower than equality of every atomic rate or convergence to
an infinite atom. At the recorded baseline, extending the maximum shell from
4 to 5 changes $n_3$ by about 7.74\% in escape and 7.58\% in Cloudy.
Both exceed the separate 5\% shell-truncation target.

\section{What MCEG manages, and why the earlier run stopped}
The new contract is \path{mceg/contracts/escape.cloudy_investigation.yaml}.
Its Equation records generate this appendix. Bindings connect named symbols
and functions to a source-pinned Subject. Retained execution files become
Evidence; scoped Assessments state what those files support. Decisions record
which approximation was selected and why. FollowUps and their evidence-backed
resolutions record the remaining work.

MCEG already supplies investigation navigation through
\texttt{mceg explain --investigation}. It did not require stopping after the
first discrepancy. The earlier provider workflow executed a finite list of
gates and retained failures, while the agent treated that implementation stage
as its stopping scope. It lacked a concrete next hypothesis, a discriminating
experiment and a completion rule requiring repeated reference agreement.
The new provider policy supplies those missing steps. The dilute-case failure
shows why the loop must continue after an apparently successful baseline repair.

Three questions remain separate:
\begin{enumerate}
\item Numerical validity: did the stated equations, conservation laws and
implicit derivatives pass their checks at the recorded source?
\item Discrepancy attribution: does an isolated intervention identify the
rate block responsible, with a reverse test and known conventions?
\item Physical adequacy: is the finite atom, transport closure and atmosphere
accurate enough for the intended observation?
\end{enumerate}
Agreement between two finite atoms answers neither infinite-atom convergence
nor planetary accuracy. The original 5\% shell-truncation criterion remains
separate and is not relaxed by this investigation.

MCEG's current FollowUp \texttt{requires} entries are human-readable conditions,
not typed executable dependencies. The provider executes cases and numerical
gates; a readiness command refuses review when a case is missing, a tolerance
fails, source identities differ, reference refinement is absent, or a required
FollowUp remains unresolved. It never creates a supported Assessment itself.
Potential MCEG improvements are a shared execution/receipt interface, optional
typed action prerequisites, and an explicit negative/refuted support status.
These would improve automation and reporting; existing records already suffice
to conduct this investigation. No MCEG runtime dependency is introduced.

\section{Evidence history and execution identity}
Old contracts, receipts and source pins remain unchanged. Native integrity
checks replay them against the historical checkout at revision
\texttt{6df616b}; they are not silently refreshed to current source. The repaired
implementation and retained-API regression receive a new Subject and Evidence.
A transient source edit/revert occurred during the first eight-case campaign.
That campaign is retained as exploratory: unchanged endpoint hashes alone
cannot prove there were no intervening edits. The final campaign must use a
frozen source identity. Replaying retained Cloudy outputs verifies their raw
hashes and creates a new escape comparison; it never reuses old escape results.
Lossless gzip storage preserves original raw hashes and receipt bytes.

The fine campaign originally labeled compressed artifacts as missing because
its loaded batch auditor preceded the gzip-aware update. Its failed batch
record remains unchanged. A separate retained audit verifies every uncompressed
hash and recomputes all eight criteria using the final auditor; no physical
calculation is rerun or result edited for that bookkeeping correction.
The retained artifacts and current MCEG Assessments are the authoritative
execution status. Passing the comparison still requires its separate numerical
and derivative checks before the complete investigation is reviewed.

\begin{thebibliography}{9}
\bibitem{badnell} Badnell et al. (2021), \emph{MNRAS} 507, 2922,
\href{https://doi.org/10.1093/mnras/stab2266}{doi:10.1093/mnras/stab2266},
equations 6--12 and 22.
\bibitem{brocklehurst} Brocklehurst (1971), \emph{MNRAS} 153, 471,
\href{https://doi.org/10.1093/mnras/153.4.471}{doi:10.1093/mnras/153.4.471},
equations 3.13 and 3.17--3.24.
\bibitem{huang} Huang et al. (2023),
\href{https://arxiv.org/html/2304.07352v1}{arXiv:2304.07352v1},
equations 12--13 for the adopted alternative $2s$--$2p$ closure.
\bibitem{cloudy} Cloudy C23.01 source,
\href{https://gitlab.nublado.org/cloudy/cloudy/-/tree/c23.01/source}{release source tree}:
\texttt{hydro\_bauman.cpp} and \texttt{CS\_l\_mixing\_PS64\_expI}.
The altered adaptations retain the Cloudy license in the package.
\end{thebibliography}
\clearpage
\appendix
\section{Generated mathematical definitions}
% Generated by mceg equations; edit the contract, then regenerate.
% Requires amsmath. Authored math is retained; TeX is not executed here.
% Definitions and hashes do not establish scientific validity.

% equation-sha256: 87d0ea395389bae387559b03d95f6c5e93c213de0443cb20a01caf527077573a
\paragraph{Equation density\_dependent\_charge\_root (revision 1)}
Density-dependent orbital rates are evaluated inside the charge-conserving root and its implicit derivative.
\begin{equation}
\begin{aligned}
P_{ij}=R_{ij}+n_e C^{e}_{ij}+n_p C^{p}_{ij}(T,n_e),\quad n_e=n_p=n_{\rm H}e^z,\\
M_{ij}=\delta_{ij}\!\left[\sum_kP_{ik}+\Gamma_i+n_eq_i^{\rm ion}\right]-P_{ji},\\
M\boldsymbol b=n_e\boldsymbol\alpha+n_e^2\boldsymbol q^{(3)},\quad n_i=n_e b_i,\\
F(z;\theta)=z+\log\!\left(1+\sum_i b_i(n_{\rm H}e^z;\theta)\right)=0,\\
\frac{\partial z}{\partial\theta}=-\frac{F_\theta}{F_z}
\end{aligned}
\label{eq:escape.cloudy_5f_investigation:density_5f_dependent_5f_charge_5f_root:1}
\end{equation}
\noindent\textbf{Scope:} Prescribed local radiation and positive photoionization; the coupled residual uses the same rate callback.

\begin{description}
\item[$T$] Common electron and proton kinetic temperature (K; ID: temperature).
\item[$n_e$] Electron density, equal to proton density in pure hydrogen (cm\textasciicircum{}-3; ID: electron\_density).
\item[$n_{\rm H}$] Total hydrogen-nucleus density (cm\textasciicircum{}-3; ID: density).
\item[$z$] Log ionized fraction (1; ID: log\_fraction).
\item[$R,C,\Gamma,\alpha,q^{\rm ion},q^{(3)}$] Radiative, binary and three-body rate coefficients (s\textasciicircum{}-1, cm\textasciicircum{}3 s\textasciicircum{}-1, cm\textasciicircum{}6 s\textasciicircum{}-1; ID: rates).
\item[$b_i$] Bound population divided by proton density (1; ID: ratios).
\item[$\theta$] External density, temperature and radiation inputs (various; ID: parameters).
\end{description}

\noindent\textbf{Assumptions:}
\begin{itemize}
\item Pure H, ne=np; the bound-state loss matrix is nonsingular.
\item The bracket is -80<=z<=0. The locally positive F\_z is checked; derivatives describe a converged root, not the bisection iterations.
\end{itemize}

% equation-sha256: 3d1b4edb827c916a9793cf31284e44806dbceaaa109656d55d4ffeec04d80e21
\paragraph{Equation psm20\_cross\_section (revision 1)}
PSM20 final-state-resolved impact-parameter cross section with both cutoffs integrated continuously.
\begin{equation}
\begin{aligned}
u(E)=\sqrt{\frac{E}{\mathrm{Ry}}\frac{m_e}{\mu}},\quad
D=\frac{3n^2\ell(n^2-\ell^2)}{2\ell+1},\quad R_1=a_0\sqrt{\frac{4D}{u^2}},\\
R_D=\sqrt{\frac{k_{\rm B}T}{4\pi n_e e^2}},\quad
R_\tau=0.72\tau\frac{a_0}{t_0}u,\quad R_c=\min(R_D,R_\tau),\\
\sigma(E)=\frac{\pi R_1^2}{4}\begin{cases}
\frac23+2\log(R_c/R_1),&R_c\ge R_1,\\
\frac23(R_c/R_1)^3,&R_c<R_1,
\end{cases}\\
q_{n\ell\to n,\ell-1}=\sqrt{\frac{8}{\pi\mu}}(k_{\rm B}T)^{-3/2}
\int_0^\infty E\sigma(E)e^{-E/(k_{\rm B}T)}\,dE
\end{aligned}
\label{eq:escape.cloudy_5f_investigation:psm20_5f_cross_5f_section:1}
\end{equation}
\noindent\textbf{Scope:} Proton-driven n=2,3 orbital mixing; same-n states are degenerate in the cross section, with actual energy differences in the inverse rate.

\begin{description}
\item[$T$] Common electron and proton kinetic temperature (K; ID: temperature).
\item[$n_e$] Electron density, equal to proton density in pure hydrogen (cm\textasciicircum{}-3; ID: electron\_density).
\item[$n$] Principal quantum number (1; ID: principal).
\item[$\ell$] Departing orbital quantum number (1; ID: orbital).
\item[$\ell-1$] Final orbital quantum number (1; ID: lower\_orbital).
\item[$\tau$] Departing state spontaneous lifetime (s; ID: lifetime).
\item[$\mu$] Neutral-protium/proton reduced mass (g; ID: mass).
\item[$E$] Collision kinetic energy (erg; ID: energy).
\item[$q$] Maxwellian binary rate coefficient (cm\textasciicircum{}3 s\textasciicircum{}-1; ID: rate).
\item[$R_1,R_D,R_\tau,R_c$] Matching, electron-Debye, lifetime and selected cutoff radii (cm; ID: radii).
\item[$a_0,t_0,\mathrm{Ry}$] Bohr length, time and Rydberg energy (cm, s, erg; ID: atomic\_units).
\end{description}

\noindent\textbf{Assumptions:}
\begin{itemize}
\item P1=1/4 for one resolved final state; the projectile charge is one.
\item The lifetime is computed from the sum of spontaneous transitions out of the departing state.
\item Unlike the C23 shortcut, the lifetime/Debye minimum is integrated at every density; this intentionally removes the shortcut jump.
\item No electron orbital-mixing channel or population fit is added to this chosen finite-atom closure.
\end{itemize}

% equation-sha256: 68d785bc96d6c29b220870e92084ae4be4e0d6e7e88a7feac9605aa66593cb99
\paragraph{Equation mixing\_inverse\_pair (revision 1)}
The paired orbital transitions obey the finite atom energy and statistical-weight convention.
\begin{equation}
\begin{aligned}
C^p_{lu}=\frac{g_u}{g_l}\exp\!\left[-\frac{E_u-E_l}{k_{\rm B}T}\right]C^p_{ul},\\
(u,l)\in\{(2p,2s),(3p,3s),(3d,3p)\},\qquad
P^{\rm mix}_{ij}=n_pC^p_{ij}
\end{aligned}
\label{eq:escape.cloudy_5f_investigation:mixing_5f_inverse_5f_pair:1}
\end{equation}
\noindent\textbf{Scope:} The finite atom replaces, rather than adds to, the previous 2s-2p fit. The historical six-state API keeps that fit.

\begin{description}
\item[$T$] Common electron and proton kinetic temperature (K; ID: temperature).
\item[$n_e$] Electron density, equal to proton density in pure hydrogen (cm\textasciicircum{}-3; ID: electron\_density).
\item[$C^p_{ij}$] Proton collision coefficient (cm\textasciicircum{}3 s\textasciicircum{}-1; ID: coefficient).
\item[$g_i$] Finite-state statistical weight (1; ID: weight).
\item[$E_i$] Finite-state excitation energy (erg; ID: excitation).
\end{description}

\noindent\textbf{Assumptions:}
\begin{itemize}
\item Pure H makes np=ne, so the paired proton rates may join the effective electron collision matrix.
\item Collapsed n=4,5 shells retain statistical orbital populations.
\end{itemize}

% equation-sha256: 31370290b36e7174c8c8b0a2348b4369761e915a926df273098add77780dcde4
\paragraph{Equation hydrogenic\_bound\_free (revision 1)}
Coulomb-wave bound-free dipoles resolve orbital states; collapsed shells average their cross sections statistically.
\begin{equation}
\begin{aligned}
t^2=\frac{h\nu}{\chi_{n\ell}}-1,\quad \kappa=t/n,\\
\sigma_{n\ell}=\frac{4\pi\alpha a_0^2n^2}{3}
\sum_{\ell'=\ell\pm1\ge0}\frac{\max(\ell,\ell')}{2\ell+1}(1+t^2)
\mathcal D_{n\ell,\ell'}^2,\\
\mathcal D_{n\ell,\ell'}^2=\frac{(n+\ell)!}{(n-\ell-1)!}(2n)^{2(\ell-n)}
\prod_{s=1}^{\ell'}(1+s^2\kappa^2)[G_\ell^{\pm}]^2,\\
\overline\sigma_n=\sum_{\ell=0}^{n-1}\frac{2\ell+1}{n^2}\sigma_{n\ell}
\end{aligned}
\label{eq:escape.cloudy_5f_investigation:hydrogenic_5f_bound_5f_free:1}
\end{equation}
\noindent\textbf{Scope:} 1<=n<=5; sigma=0 below the actual finite-state threshold. The radial functions are defined by the following recurrence.

\begin{description}
\item[$n$] Principal quantum number (1; ID: principal).
\item[$\ell$] Bound-state orbital quantum number (1; ID: orbital).
\item[$h\nu/\chi_{n\ell}$] Photon energy divided by the level threshold (1; ID: relative\_energy).
\item[$\nu$] Photon frequency (Hz; ID: frequency).
\item[$\sigma_{n\ell}$] Orbital-resolved photoionization cross section (cm\textasciicircum{}2; ID: cross\_section).
\item[$\mathcal D,G^\pm$] Dimensionless Coulomb radial factors (1; ID: radial).
\item[$\chi_{n\ell}$] Finite-level ionization threshold (erg; ID: threshold).
\end{description}

\noindent\textbf{Assumptions:}
\begin{itemize}
\item Resolved states are 1s,2s,2p,3s,3p,3d; n=4,5 are collapsed.
\item Hydrogenic Coulomb dipoles are evaluated at the selected level threshold.
\item No fitted correction is applied to n2/n3 populations.
\end{itemize}

% equation-sha256: 9861b547e50970f84ebd2d0f5827714e37bcb847a1ad02bfc057b015cda2eddf
\paragraph{Equation brocklehurst\_recurrence (revision 1)}
Finite downward recurrence used to evaluate the hydrogenic radial factors without population fitting.
\begin{equation}
\begin{aligned}
G_0^*=\sqrt{\frac\pi2}\frac{8n(4n)^n}{(2n-1)!}
\frac{e^{-2n\arctan(t)/t}}{\sqrt{1-e^{-2\pi/\kappa}}(1+t^2)^{n+2}},\\
G^+_{n-1}=G_0^*,\quad G^+_{n-2}=n(2n-1)(1+t^2)G_0^*,\\
G^-_{n-1}=\frac{1+t^2}{2n}G_0^*,\quad
G^-_{n-2}=(2n-1)[4+(n-1)(1+t^2)]G^-_{n-1},\\
A_q^+=4n^2-4(q+2)^2+(q+2)(2q+3)(1+t^2),\\
B_q^+=4n^2[n^2-(q+2)^2][1+(q+3)^2\kappa^2],\\
A_q^-=4n^2-4(q+1)^2+(q+1)(2q+3)(1+t^2),\\
B_q^-=4n^2[n^2-(q+2)^2][1+(q+1)^2\kappa^2],\\
G_q^\pm=A_q^\pm G_{q+1}^\pm-B_q^\pm G_{q+2}^\pm
\end{aligned}
\label{eq:escape.cloudy_5f_investigation:brocklehurst_5f_recurrence:1}
\end{equation}
\noindent\textbf{Scope:} The plus recurrence runs q=n-3,...,0 and the minus recurrence q=n-3,...,1. Use initial values only for existing indices.

\begin{description}
\item[$n$] Principal quantum number (1; ID: principal).
\item[$\ell$] Bound-state orbital quantum number (1; ID: orbital).
\item[$1+t^2$] Photon energy divided by threshold (1; ID: relative\_energy).
\item[$G_q^\pm$] Radial recurrence quantities (1; ID: radial).
\item[$\kappa=t/n$] Dimensionless continuum momentum (1; ID: continuum).
\end{description}

\noindent\textbf{Assumptions:}
\begin{itemize}
\item At threshold arctan(t)/t approaches one and exp(-2pi/kappa) approaches zero.
\item Small-t series avoids a singular intermediate and preserves differentiation.
\item Ordinary float64 is restricted to n<=5; high-n overflow behavior is not inferred.
\end{itemize}

% equation-sha256: 9e5010ac050dd8f0a805455aa753cc2eea3484186b90bea6cf1d0677c49d9024
\paragraph{Equation paired\_continuum\_kernels (revision 1)}
One photoionization cross section supplies its inverse Milne rate and continuum absorption/emission.
\begin{equation}
\begin{aligned}
\eta_\nu=\frac{c^2J_\nu}{2h\nu^3},\quad
 a_i=\frac{g_i}{2}\left(\frac{h^2}{2\pi m_e k_{\rm B}T}\right)^{3/2}
 e^{-(h\nu-\chi_i)/(k_{\rm B}T)},\\
\Gamma_i=\frac{8\pi}{c^2}\int\nu^2\sigma_i\eta_\nu\,d\nu,\quad
\alpha_i=\frac{8\pi}{c^2}\int\nu^2\sigma_i a_i(1+\eta_\nu)\,d\nu,\\
\chi_\nu^{\rm bf}=\sum_i\sigma_i(n_i-n_en_pa_i),\quad
j_\nu^{\rm bf}=\frac{2h\nu^3}{c^2}n_en_p\sum_i\sigma_i a_i
\end{aligned}
\label{eq:escape.cloudy_5f_investigation:paired_5f_continuum_5f_kernels:1}
\end{equation}
\noindent\textbf{Scope:} Included bound-free processes only; other free-free and line kernels are unchanged and separately recorded in the earlier coupled formulation.

\begin{description}
\item[$T$] Common electron and proton kinetic temperature (K; ID: temperature).
\item[$n_e$] Electron density, equal to proton density in pure hydrogen (cm\textasciicircum{}-3; ID: electron\_density).
\item[$\nu$] Frequency (Hz; ID: frequency).
\item[$\eta_\nu$] Photon occupation number (1; ID: occupation).
\item[$d\nu$] Frequency quadrature measure (Hz; ID: weight).
\item[$n_i$] Bound population (cm\textasciicircum{}-3; ID: population).
\item[$\sigma_i$] The same orbital-resolved or collapsed cross section in every kernel (cm\textasciicircum{}2; ID: cross\_section).
\item[$\Gamma_i,\alpha_i$] Photoionization and radiative recombination coefficients (s\textasciicircum{}-1, cm\textasciicircum{}3 s\textasciicircum{}-1; ID: rates).
\item[$\chi_\nu^{\rm bf},j_\nu^{\rm bf}$] Bound-free extinction and emissivity (cm\textasciicircum{}-1, erg s\textasciicircum{}-1 cm\textasciicircum{}-3 sr\textasciicircum{}-1 Hz\textasciicircum{}-1; ID: transfer).
\end{description}

\noindent\textbf{Assumptions:}
\begin{itemize}
\item Maxwellian electrons at the gas temperature; stimulated recombination is retained.
\item The same selected quadrature and cross section must be used in local rates and coupled continuum transfer.
\end{itemize}

\end{document}
